
\subsubsection*{3.2.1}
\phantomsection
\label{sga3:VIIB:3.2.1}

The bijection \(\exp_{U(L(C))}\) allows us, by transport of
structure, to define a group law on the set
\[
  L(C)\otimes_C\mathfrak r(C),
\]
which, as in \SGARefLink{sga3:VIIB:3.2:prop}{3.2}, we identify with a subset of \(U(L(C))\). If
\(x\) and \(y\) are two elements of
\(L(C)\otimes_C\mathfrak r(C)\), this law is given by
\phantomsection
\label{sga3:VIIB:3.2.1:transported-product}
\[
\begin{aligned}
  x\mathbin{\cdot}y
  &=
  \log\bigl((\exp x)(\exp y)\bigr)
  =
  \log\left(
    1+\sum_{p+q>0}\frac{x^p y^q}{p!\,q!}
  \right)
  \\
  &=
  \sum_{m\geq1}
  \sum_{p_i+q_i>0}
  \frac{(-1)^{m-1}}{m}
  \frac{x^{p_1}y^{q_1}}{p_1!\,q_1!}
  \cdots
  \frac{x^{p_m}y^{q_m}}{p_m!\,q_m!}
  \\
  &=
  \sum_{\ell\geq1}P_\ell(x,y).
\end{aligned}
\]
Here \(P_\ell(x,y)\) denotes the sum of the monomials of total
degree \(\ell\) in \(x\) and \(y\). For example,
\[
  P_1(x,y)=x+y,
\]
and
\[
\begin{aligned}
  P_2(x,y)
  &=
  \underbrace{
    \frac{x^2}{2}+xy+\frac{y^2}{2}
  }_{m=1}
  -
  \underbrace{
    \frac12\bigl(x^2+xy+yx+y^2\bigr)
  }_{m=2}
  \\
  &=
  \frac12(xy-yx)
  =
  \frac12[x,y].
\end{aligned}
\]
Moreover, \(P_3(x,y)\) is the sum of the following three terms:
\[
\begin{aligned}
  P_3(x,y)
  ={}&
  \underbrace{
    \frac{x^3}{6}
    +\frac{x^2y}{2}
    +\frac{xy^2}{2}
    +\frac{y^3}{6}
  }_{m=1}
  \\
  &-
  \underbrace{
    \frac12\left(
      x^3+\frac32x^2y+\frac12yx^2+xyx+yxy
      +\frac12y^2x+\frac32xy^2+y^3
    \right)
  }_{m=2}
  \\
  &+
  \underbrace{
    \frac13\left(
      x^3+x^2y+yx^2+xyx+yxy+y^2x+xy^2+y^3
    \right)
  }_{m=3}.
\end{aligned}
\]
Consequently,
\phantomsection
\label{sga3:VIIB:3.2.1:P3}
\[
  P_3(x,y)
  =
  \frac1{12}
  \bigl(
    x^2y+yx^2-2xyx-2yxy+y^2x+xy^2
  \bigr)
  =
  \frac1{12}
  \bigl(
    [[y,x],x]+[y,[y,x]]
  \bigr).
\]

More generally, one can prove the Campbell--Hausdorff
formula:
\phantomsection
\label{sga3:VIIB:3.2.1:campbell-hausdorff}
\[
\begin{aligned}
  P_\ell(x,y)
  ={}&
  \sum_{m=1}^{\ell}
  \frac{(-1)^{m-1}}{m\cdot\ell}
  \sum_{\substack{
    p_1,\ldots,p_m\\
    q_1,\ldots,q_{m-1}
  }}
  \left(
    \prod_{i=1}^{m-1}
    \frac{(\operatorname{ad}x)^{p_i}}{p_i!}
    \frac{(\operatorname{ad}y)^{q_i}}{q_i!}
  \right)
  \frac{(\operatorname{ad}x)^{p_m}}{p_m!}(y)
  \\
  &+
  \sum_{m=1}^{\ell}
  \frac{(-1)^{m-1}}{m\cdot\ell}
  \sum_{\substack{
    p_1,\ldots,p_{m-1}\\
    q_1,\ldots,q_{m-1}
  }}
  \left(
    \prod_{i=1}^{m-1}
    \frac{(\operatorname{ad}x)^{p_i}}{p_i!}
    \frac{(\operatorname{ad}y)^{q_i}}{q_i!}
  \right)(x).
\end{aligned}
\]
Here the \(p_j,q_i\in\mathbb N\) satisfy
\[
  p_i+q_i\geq1
  \quad(1\leq i\leq m-1),
  \qquad
  p_m+\sum_{i=1}^{m-1}(p_i+q_i)=\ell-1.
\]
In other words, in the sums above every nonzero ``Lie monomial''
has total degree \(\ell\). For a proof, see N.~Jacobson,
\emph{Lie Algebras} (Interscience, 1962), \S V.5, or
\textup{[BLie]}, II, \S6.4, Theorem~2.
